\clearpage
\begin{examplebox}
\par\medskip
\noindent\begin{minipage}{\linewidth}
\paragraph{例1：两个零点之间，直接找到导数等于零的位置}
\begin{center}% F081: precise mathematical teaching diagram.
\begin{tikzpicture}[x=1cm,y=1cm,>=stealth,font=\small]
\draw[rounded corners=3pt,gray!40] (0,0) rectangle (14.2,9.0);
\begin{scope}[yshift=1cm]
\node at (7.1,7.4) {例1：两个零点之间找到水平切线};
\begin{scope}[shift={(5.4,6)},x=1.55cm,y=0.8cm]
\draw[->,gray!65] (-2.4,0)--(2.45,0) node[right] {$x$};
\draw[->,gray!65] (0,-4.8)--(0,0.7) node[above] {$y$};
\draw[blue!75!black,line width=1.15pt,,domain=-2:2,samples=160] plot (\x,{(\x)*(\x)-4});
\fill[blue!75!black] (-2,0) circle (1.8pt);\node[above left] at (-2,0) {$a=-2$};
\fill[blue!75!black] (2,0) circle (1.8pt);\node[above right] at (2,0) {$b=2$};
\fill[blue!75!black] (0,-4) circle (1.8pt);\node[below right] at (0,-4) {$c=0$};
\draw[orange!85!black,line width=1.2pt] (-0.65,-4)--(0.65,-4);
\end{scope}
\node[] at (10.8,3.5) {$f^{\prime}(c)=0$};
\end{scope}
\node[text width=13cm,align=center] at (7.1,.9) {$f(x)=x^2-4$，两零点 $-2,2$ 夹着 $c=0$。};
\end{tikzpicture}
\end{center}
\end{minipage}\par\medskip

已知
\[
f(x)=x^2-4.
\]
\textbf{先找零点：}解 $x^2-4=0$，得 $x=-2$ 或 $x=2$。

\textbf{再用罗尔定理：}函数图像连续，区间内部处处可导，又有
\[
f(-2)=f(2)=0,
\]
所以在 $(-2,2)$ 内至少有一个 $c$ 使 $f'(c)=0$。

\textbf{直接算一遍：}
\[
f'(x)=2x,\qquad f'(0)=0.
\]
因此可以取 $c=0$。

\textbf{一道题记住：}有两个不同零点 $\Longrightarrow$ 中间至少一个导数等于零的位置；但这个位置不一定是两个零点的正中间。

\end{examplebox}
\clearpage
\begin{examplebox}
\par\medskip
\noindent\begin{minipage}{\linewidth}
\paragraph{例2：根本不用解三次方程，就能判断恰好三个零点}
\begin{center}% F081: precise mathematical teaching diagram.
\begin{tikzpicture}[x=1cm,y=1cm,>=stealth,font=\small]
\draw[rounded corners=3pt,gray!40] (0,0) rectangle (14.2,9.0);
\begin{scope}[yshift=1cm]
\node at (7.1,7.4) {例2：两个导数零点限定上限，三段异号区间保证存在};
\begin{scope}[shift={(6.1,3.1)},x=2cm,y=0.85cm]
\draw[->,gray!65] (-2.35,0)--(2.4,0) node[right] {$x$};
\draw[->,gray!65] (0,-1.8)--(0,3.9) node[above] {$y$};
\draw[blue!75!black,line width=1.15pt,,domain=-2.1:2.1,samples=160] plot (\x,{(\x)*(\x)*(\x)-3*\x+1});
\fill[blue!75!black] (-2,-1) circle (1.8pt);\node[above left] at (-2,-1) {$(-2,-1)$};
\fill[blue!75!black] (-1,3) circle (1.8pt);\node[above left] at (-1,3) {$(-1,3)$};
\fill[blue!75!black] (1,-1) circle (1.8pt);\node[below right] at (1,-1) {$(1,-1)$};
\fill[blue!75!black] (2,3) circle (1.8pt);\node[above left] at (2,3) {$(2,3)$};
\fill[green!50!black] (-1.8793852415718,0) circle (1.8pt);
\fill[green!50!black] (0.3472963553339,0) circle (1.8pt);
\fill[green!50!black] (1.532088886238,0) circle (1.8pt);
\draw[orange!85!black,line width=1.2pt] (-1.35,3)--(-0.65,3);
\draw[orange!85!black,line width=1.2pt] (0.65,-1)--(1.35,-1);
\end{scope}
\node[] at (7.1,1) {三段异号区间 $(-2,-1),(-1,1),(1,2)$ 各至少一根。};
\end{scope}
\node[text width=13cm,align=center] at (7.1,.9) {$f=x^3-3x+1$，$f^{\prime}$ 只有两个零点，所以原函数恰好三个不同零点。};
\end{tikzpicture}
\end{center}
\end{minipage}\par\medskip

判断
\[
f(x)=x^3-3x+1
\]
有几个不同的实数零点。

\textbf{第一步，先求最多有几个。}
\[
f'(x)=3(x+1)(x-1).
\]
导数等于 $0$ 的位置只有 $x=-1$、$x=1$ 两个，所以根据罗尔定理，$f(x)$ \textbf{最多只有三个不同零点}。

\textbf{第二步，再证明至少有三个。}计算：
\[
\begin{array}{c|rrrr}
x&-2&-1&1&2\\ \hline
f(x)&-1&3&-1&3
\end{array}
\]
从 $-2$ 到 $-1$，函数值由负变正；从 $-1$ 到 $1$，由正变负；从 $1$ 到 $2$，由负变正。因为函数图像连续，所以这三个小区间内\textbf{各至少有一个零点}。

因此\textbf{至少三个、最多三个}，得到
\[
\boxed{\ f(x)=0\text{ 恰好有三个不同实数解}\ }.
\]
\textbf{最值得学的思路：}\textbf{先用导数管住上限，再用函数值的正负变化凑足数量。}
\end{examplebox}
\clearpage
\begin{examplebox}
\par\medskip
\noindent\begin{minipage}{\linewidth}
\paragraph{例3：导数没有零点，原方程最多一个解}
\begin{center}% F081: precise mathematical teaching diagram.
\begin{tikzpicture}[x=1cm,y=1cm,>=stealth,font=\small]
\draw[rounded corners=3pt,gray!40] (0,0) rectangle (14.2,8.6);
\begin{scope}[yshift=1cm]
\node at (7.1,7.0) {例3：上限为一个，再用异号端点保证存在};
\begin{scope}[shift={(4.5,3.3)},x=7cm,y=1.0cm]
\draw[->,gray!65] (-0.08,0)--(0.36,0) node[right] {$x-1$};
\draw[->,gray!65] (0,-1.8)--(0,2.4) node[above] {$y$};
\draw[blue!75!black,line width=1.15pt,,domain=-0.06:0.28,samples=160] plot (\x,{(\x+1)*(\x+1)*(\x+1)*(\x+1)*(\x+1)+(\x+1)-3});
\fill[blue!75!black] (0,-1) circle (1.8pt);\node[above left] at (0,-1) {$(1,-1)$};
\end{scope}
\node[] at (10.5,5.5) {$f^{\prime}(x)=5x^4+1>0$};
\node[] at (10.5,4.4) {全实数范围最多一根};
\node[] at (10.5,3.2) {$f(1)=-1,\quad f(2)=31$};
\node[] at (10.5,2.1) {连续保证至少一根};
\end{scope}
\node[text width=13cm,align=center] at (7.1,.9) {左图放大交点附近，横轴为 $x-1$；唯一零点位于 $(1,2)$。};
\end{tikzpicture}
\end{center}
\end{minipage}\par\medskip

判断
\[
f(x)=x^5+x-3
\]
有几个实数零点。

\textbf{先看导数：}
\[
f'(x)=5x^4+1>0.
\]
因为导数无论在什么位置都不等于 $0$，所以根据罗尔定理，原函数在整个实数范围内\textbf{不可能有两个不同零点}。

\textbf{再找是否有解：}
\[
f(1)=-1<0,\qquad f(2)=31>0.
\]
函数图像连续，$x=1$ 到 $x=2$ 之间一定会与 $x$ 轴相交一次。

所以
\[
\boxed{\ f(x)=0\text{ 恰好有一个实数解，且在 }(1,2)\text{ 内}\ }.
\]
\textbf{小提示：}这里的“最多一个”，既可以从导数始终大于 $0$ 看出，也可以由罗尔定理直接得出。

\end{examplebox}
\clearpage
\begin{examplebox}
\par\medskip
\noindent\begin{minipage}{\linewidth}
\paragraph{例4：有参数时，不用求三次方程的通用解法}
\begin{center}% F081: precise mathematical teaching diagram.
\begin{tikzpicture}[x=1cm,y=1cm,>=stealth,font=\small]
\draw[rounded corners=3pt,gray!40] (0,0) rectangle (14.2,9.2);
\begin{scope}[yshift=1cm]
\node at (7.1,7.6) {例4：只改变参数 $t$，曲线整体上下移动};
\node[] at (2.6,7) {$t=0$：3个不同零点};
\begin{scope}[shift={(2.6,3.8)},x=0.54cm,y=0.38cm]
\draw[->,gray!65] (-3,0)--(3,0) node[right] {$x$};
\draw[->,gray!65] (0,-6)--(0,6) node[above] {$y$};
\begin{scope}\clip (-2.8,-5.5) rectangle (2.8,5.5);
\draw[blue!75!black,line width=1.15pt,,domain=-2.8:2.8,samples=160] plot (\x,{(\x)*(\x)*(\x)-3*\x+(0)});
\end{scope}
\fill[blue!75!black] (-1.7320508075689,0) circle (1.8pt);
\fill[blue!75!black] (-0.0,0) circle (1.8pt);
\fill[blue!75!black] (1.7320508075689,0) circle (1.8pt);
\fill[orange!85!black] (-1,2) circle (1.8pt);
\draw[orange!85!black,line width=1.2pt] (-1.32,2)--(-0.6799999999999999,2);
\fill[orange!85!black] (1,-2) circle (1.8pt);
\draw[orange!85!black,line width=1.2pt] (0.6799999999999999,-2)--(1.32,-2);
\end{scope}
\node[] at (2.6,1.15) {$f(-1)=2,\ f(1)=-2$};
\node[] at (7.1,7) {$t=2$：2个不同零点};
\begin{scope}[shift={(7.1,3.8)},x=0.54cm,y=0.38cm]
\draw[->,gray!65] (-3,0)--(3,0) node[right] {$x$};
\draw[->,gray!65] (0,-6)--(0,6) node[above] {$y$};
\begin{scope}\clip (-2.8,-5.5) rectangle (2.8,5.5);
\draw[blue!75!black,line width=1.15pt,,domain=-2.8:2.8,samples=160] plot (\x,{(\x)*(\x)*(\x)-3*\x+(2)});
\end{scope}
\fill[blue!75!black] (-2.0,0) circle (1.8pt);
\fill[blue!75!black] (1,0) circle (1.8pt);
\fill[orange!85!black] (-1,4) circle (1.8pt);
\draw[orange!85!black,line width=1.2pt] (-1.32,4)--(-0.6799999999999999,4);
\fill[orange!85!black] (1,0) circle (1.8pt);
\draw[orange!85!black,line width=1.2pt] (0.6799999999999999,0)--(1.32,0);
\end{scope}
\node[] at (7.1,1.15) {$f(-1)=4,\ f(1)=0$};
\node[] at (11.6,7) {$t=3$：1个不同零点};
\begin{scope}[shift={(11.6,3.8)},x=0.54cm,y=0.38cm]
\draw[->,gray!65] (-3,0)--(3,0) node[right] {$x$};
\draw[->,gray!65] (0,-6)--(0,6) node[above] {$y$};
\begin{scope}\clip (-2.8,-5.5) rectangle (2.8,5.5);
\draw[blue!75!black,line width=1.15pt,,domain=-2.8:2.8,samples=160] plot (\x,{(\x)*(\x)*(\x)-3*\x+(3)});
\end{scope}
\fill[blue!75!black] (-2.1038034027355,0) circle (1.8pt);
\fill[orange!85!black] (-1,5) circle (1.8pt);
\draw[orange!85!black,line width=1.2pt] (-1.32,5)--(-0.6799999999999999,5);
\fill[orange!85!black] (1,1) circle (1.8pt);
\draw[orange!85!black,line width=1.2pt] (0.6799999999999999,1)--(1.32,1);
\end{scope}
\node[] at (11.6,1.15) {$f(-1)=5,\ f(1)=1$};
\end{scope}
\node[text width=13cm,align=center] at (7.1,.9) {各面板同坐标尺度；横轴上的点按不同位置计数，相切也算一个零点。};
\end{tikzpicture}
\end{center}
\end{minipage}\par\medskip

求方程
\[
x^3-3x+t=0\qquad (t\text{ 为实数})
\]
的\textbf{不同实数解个数}。

令 $f(x)=x^3-3x+t$，则
\[
f'(x)=3(x+1)(x-1).
\]
无论 $t$ 取什么值，导数都只在 $-1,1$ 两处等于 $0$，所以\textbf{最多三个不同解}。

再算出两个容易算的位置的函数值：
\[
f(-1)=t+2,\qquad f(1)=t-2.
\]
导数的正负号说明函数先升、再降、再升；当 $x$ 足够小时函数值为负，足够大时为正。因此：
\[
\boxed{
\begin{array}{c|c}
t\text{ 的范围}&\text{不同实数解个数}\\ \hline
-2<t<2&3\\
t=-2\text{ 或 }t=2&2\\
t<-2\text{ 或 }t>2&1
\end{array}}
\]
\begin{center}% F081: precise mathematical teaching diagram.
\begin{tikzpicture}[x=1cm,y=1cm,>=stealth,font=\small]
\draw[rounded corners=3pt,gray!40] (0,0) rectangle (14.2,9.2);
\begin{scope}[yshift=1cm]
\node at (7.1,7.6) {例4：只改变参数 $t$，曲线整体上下移动};
\node[] at (2.6,7) {$t=0$：3个不同零点};
\begin{scope}[shift={(2.6,3.8)},x=0.54cm,y=0.38cm]
\draw[->,gray!65] (-3,0)--(3,0) node[right] {$x$};
\draw[->,gray!65] (0,-6)--(0,6) node[above] {$y$};
\begin{scope}\clip (-2.8,-5.5) rectangle (2.8,5.5);
\draw[blue!75!black,line width=1.15pt,,domain=-2.8:2.8,samples=160] plot (\x,{(\x)*(\x)*(\x)-3*\x+(0)});
\end{scope}
\fill[blue!75!black] (-1.7320508075689,0) circle (1.8pt);
\fill[blue!75!black] (-0.0,0) circle (1.8pt);
\fill[blue!75!black] (1.7320508075689,0) circle (1.8pt);
\fill[orange!85!black] (-1,2) circle (1.8pt);
\draw[orange!85!black,line width=1.2pt] (-1.32,2)--(-0.6799999999999999,2);
\fill[orange!85!black] (1,-2) circle (1.8pt);
\draw[orange!85!black,line width=1.2pt] (0.6799999999999999,-2)--(1.32,-2);
\end{scope}
\node[] at (2.6,1.15) {$f(-1)=2,\ f(1)=-2$};
\node[] at (7.1,7) {$t=-2$：2个不同零点};
\begin{scope}[shift={(7.1,3.8)},x=0.54cm,y=0.38cm]
\draw[->,gray!65] (-3,0)--(3,0) node[right] {$x$};
\draw[->,gray!65] (0,-6)--(0,6) node[above] {$y$};
\begin{scope}\clip (-2.8,-5.5) rectangle (2.8,5.5);
\draw[blue!75!black,line width=1.15pt,,domain=-2.8:2.8,samples=160] plot (\x,{(\x)*(\x)*(\x)-3*\x+(-2)});
\end{scope}
\fill[blue!75!black] (-1,0) circle (1.8pt);
\fill[blue!75!black] (2.0,0) circle (1.8pt);
\fill[orange!85!black] (-1,0) circle (1.8pt);
\draw[orange!85!black,line width=1.2pt] (-1.32,0)--(-0.6799999999999999,0);
\fill[orange!85!black] (1,-4) circle (1.8pt);
\draw[orange!85!black,line width=1.2pt] (0.6799999999999999,-4)--(1.32,-4);
\end{scope}
\node[] at (7.1,1.15) {$f(-1)=0,\ f(1)=-4$};
\node[] at (11.6,7) {$t=-3$：1个不同零点};
\begin{scope}[shift={(11.6,3.8)},x=0.54cm,y=0.38cm]
\draw[->,gray!65] (-3,0)--(3,0) node[right] {$x$};
\draw[->,gray!65] (0,-6)--(0,6) node[above] {$y$};
\begin{scope}\clip (-2.8,-5.5) rectangle (2.8,5.5);
\draw[blue!75!black,line width=1.15pt,,domain=-2.8:2.8,samples=160] plot (\x,{(\x)*(\x)*(\x)-3*\x+(-3)});
\end{scope}
\fill[blue!75!black] (2.1038034027355,0) circle (1.8pt);
\fill[orange!85!black] (-1,-1) circle (1.8pt);
\draw[orange!85!black,line width=1.2pt] (-1.32,-1)--(-0.6799999999999999,-1);
\fill[orange!85!black] (1,-5) circle (1.8pt);
\draw[orange!85!black,line width=1.2pt] (0.6799999999999999,-5)--(1.32,-5);
\end{scope}
\node[] at (11.6,1.15) {$f(-1)=-1,\ f(1)=-5$};
\end{scope}
\node[text width=13cm,align=center] at (7.1,.9) {各面板同坐标尺度；横轴上的点按不同位置计数，相切也算一个零点。};
\end{tikzpicture}
\end{center}

\textbf{等号为什么是两个？}例如 $t=2$ 时
\[
x^3-3x+2=(x-1)^2(x+2),
\]
只有 $x=1$ 和 $x=-2$ 两个\textbf{不同}解；$t=-2$ 时同理，只有 $x=-1$ 和 $x=2$。

\textbf{这一题的亮点：}求导后导数里没有 $t$，意味着不论参数怎样变化，\textbf{零点的最多数量始终是三个}；参数只改变实际出现一个、两个还是三个。
\end{examplebox}
\clearpage
\begin{examplebox}
\par\medskip
\noindent\begin{minipage}{\linewidth}
\paragraph{例5：连续用两次罗尔定理，可以省下大量计算}
\begin{center}% F081: precise mathematical teaching diagram.
\begin{tikzpicture}[x=1cm,y=1cm,>=stealth,font=\small]
\draw[rounded corners=3pt,gray!40] (0,0) rectangle (14.2,8.6);
\begin{scope}[yshift=1cm]
\node at (7.1,7.0) {例5：假设三个零点，就会逼出一个二阶导数零点};
\draw[->,gray!70] (3,5.9)--(12.3,5.9);
\node[] at (1.55,5.9) {$f=0$};
\fill[blue!75!black] (3.5999999999999996,5.9) circle (2pt);\node[] at (3.5999999999999996,6.300000000000001) {$a$};
\fill[blue!75!black] (7.1,5.9) circle (2pt);\node[] at (7.1,6.300000000000001) {$b$};
\fill[blue!75!black] (10.6,5.9) circle (2pt);\node[] at (10.6,6.300000000000001) {$c$};
\draw[->,gray!70] (3,3.9)--(12.3,3.9);
\node[] at (1.55,3.9) {$f^{\prime}=0$};
\fill[blue!75!black] (5.07927405783631,3.9) circle (2pt);\node[] at (5.07927405783631,4.3) {$u$};
\fill[blue!75!black] (9.12072594216369,3.9) circle (2pt);\node[] at (9.12072594216369,4.3) {$v$};
\draw[->,gray!70] (3,1.9)--(12.3,1.9);
\node[] at (1.55,1.9) {$f^{\prime\prime}=0$};
\fill[blue!75!black] (7.1,1.9) circle (2pt);\node[] at (7.1,2.3) {$w$};
\draw[dashed,orange!75!black] (3.5999999999999996,5.75)--(5.07927405783631,4.55);
\draw[dashed,orange!75!black] (7.1,5.75)--(5.07927405783631,4.55);
\draw[dashed,orange!75!black] (7.1,5.75)--(9.12072594216369,4.55);
\draw[dashed,orange!75!black] (10.6,5.75)--(9.12072594216369,4.55);
\draw[dashed,orange!75!black] (5.07927405783631,3.75)--(7.1,2.55);
\draw[dashed,orange!75!black] (9.12072594216369,3.75)--(7.1,2.55);
\end{scope}
\node[text width=13cm,align=center] at (7.1,.9) {假设所得的位置排列（示意）；末行 $f^{\prime\prime}(w)=0$ 与 $f^{\prime\prime}>0$ 矛盾。};
\end{tikzpicture}
\end{center}
\end{minipage}\par\medskip

已知 $f$ 在同一个区间 $I$ 上连续，在其内部二阶可导，且内部处处有
\[
f''(x)>0.
\]
问 $f(x)=0$ 最多有几个不同解？

\textbf{假设有三个：}记为 $a<b<c$。

由罗尔定理，在 $(a,b)$ 和 $(b,c)$ 内至少各有一个 $f'(x)=0$ 的位置，记为 $u<v$。

再对 $f'(x)$ 在 $[u,v]$ 上用一次罗尔定理，就得到某个 $w\in(u,v)$ 满足
\[
f''(w)=0.
\]
这和 $f''(x)>0$ 矛盾。

所以
\[
\boxed{\ f(x)=0\text{ 最多只有两个不同解}\ }.
\]
\textbf{记忆方法：}原函数有三个零点 $\Rightarrow$ 第一次求导至少两个 $0$ $\Rightarrow$ 第二次求导至少一个 $0$。如果最后那个 $0$ 根本不可能出现，最初也就不可能有三个零点。
\end{examplebox}

\clearpage
\begin{examplebox}
\paragraph{例6：两点不在 $x$ 轴上，罗尔定理照样成立}

设
\[
 f(x)=(x-2)^2+1.
\]
\textbf{第一步，找出两个同高的位置：}
\[
 f(1)=2,\qquad f(3)=2.
\]
于是曲线上有 $A(1,2)$ 和 $B(3,2)$，直线 $AB$ 水平，其斜率为 $0$。

\textbf{第二步，找中间与 $AB$ 平行的切线：}
\[
 f'(x)=2(x-2),\qquad f'(2)=0.
\]
所以切点是 $C(2,1)$，其切线为 $y=1$，与 $AB$ 所在的 $y=2$ \textbf{平行}。

\textbf{第三步，理解两层意思：}虽然 $C$ 比 $A,B$ 低，罗尔定理仍然成立；而把 $y=2$ 看成“横向划过图像的一条线”，它与图像有两个交点，中间就必定夹着至少一处水平切线。

\[
\boxed{\ \text{两点同高}\ \Longrightarrow\ \text{中间存在与连线平行的水平切线}\ }
\]
\textbf{提示：}如果把 $y=2$ 改为任意 $y=h$，求交点个数仍可以使用“先数水平切线位置，再限制交点数量”的方法。
\end{examplebox}


% -------- 新增例题：一点导数为零，不同走向 --------
\clearpage
\begin{examplebox}
\paragraph{例7：同样在原点切线水平，为什么有的掉头、有的不断向前？}

四个函数在 $x=0$ 都满足 $f'(0)=0$，但左右走向不同：
\[
\begin{array}{c|c|c|c}
\text{函数}&x<0\text{ 时的导数}&x>0\text{ 时的导数}&\text{结果}\\ \hline
-x^2&-2x>0&-2x<0&\text{先升后降（最高处）}\\
\phantom{-}x^2&2x<0&2x>0&\text{先降后升（最低处）}\\
\phantom{-}x^3&3x^2>0&3x^2>0&\text{一直上升（不掉头）}\\
-x^3&-3x^2<0&-3x^2<0&\text{一直下降（不掉头）}
\end{array}
\]

\textbf{最快判断：}先算 $f'(0)$，再挑 $0$ 左边与右边各一个小范围，
分别看导数是正还是负；\textbf{正变负才是从升到降，负变正才是从降到升}。

\textbf{联系罗尔定理：}两端同高可以确保中间找到水平切线；
但如果只给你一个 $f'(c)=0$，\textbf{不能反过来认定两端同高，也不能断言一定掉头}。
\end{examplebox}
